\bookchapter{제5장: 스마트폰이 나를 더 잘 이해하도록 - 기본 설정}

\chapterintro{몇 분만 설정하면 사용성이 크게 향상됩니다. 이 설정들은 스마트폰을 “조율”하는 것과 같아서, 사용 습관에 더 잘 맞게 만들어 줍니다.}

\booksection{사운드와 진동}

마음에 드는 벨소리와 알림음을 설정하세요. 무음 모드(회의), 방해 금지 모드(수면)를 유연하게 사용하면 세상이 순식간에 고요해집니다.

\begin{itemize}
  \item \textbf{벨소리}: 들을 수 있고, 어색하지 않은 것을 선택하세요.
  \item \textbf{진동 모드}: 무음일 때 진동으로 알림.
  \item \textbf{방해 금지 모드}: 예외를 설정할 수 있으며, 예: 주소록의 전화는 허용.
  \item \textbf{볼륨 분리}: 미디어 볼륨, 벨소리 볼륨, 알람 볼륨을 별도로 조절할 수 있습니다.
\end{itemize}

\booksection{디스플레이와 밝기}

\begin{itemize}
  \item \textbf{자동 밝기}: 스마트폰이 환경광에 따라 자동 조절, 마음 편함.
  \item \textbf{눈 보호/나이트 모드}: 블루라이트를 필터링하여 밤에 스마트폰을 볼 때 편안합니다.
  \item \textbf{다크 모드}: 전력 절약(OLED 화면에 특히 효과적)되며, 야간에 눈이 부시지 않습니다. 많은 앱도 다크 모드를 지원합니다.
  \item \textbf{화면 주사율}: 일부 안드로이드 기기는 조절할 수 있으며, 높을수록(예: 120Hz) 스크롤이 더 부드러워집니다.
  \item \textbf{자동 잠금}: 15초 또는 30초로 설정하여 안전하고 전력 절약됩니다.
  \item \textbf{글꼴 크기}: 글자가 작다고 느끼면 크게 조절하세요.
\end{itemize}

\booksection{배경화면, 테마 및 데스크톱}

예쁜 배경화면(정지 또는 동적)으로 바꾸면 기분도 달라집니다. 안드로이드 브랜드 스마트폰에는 종종 테마 스토어가 있어, 아이콘과 스타일을 한꺼번에 변경할 수 있습니다.

\textbf{데스크톱 위젯}: 날씨, 일정, 음악 제어 등을 데스크톱에 배치하여 한 번에 접근할 수 있습니다. iOS와 안드로이드 모두 지원하며, 데스크톱 빈 공간을 길게 눌러 추가합니다.

\booksection{효율성 설정: 스마트폰을 “외장 두뇌”로 만들기}

이 부분은 “한번 설정하면 못 돌아갈” 기능들로, 각각 일상 효율성을 크게 향상시킬 수 있습니다.

\booksubsection{집중 / 방해금지 모드: 알림에 빼앗긴 집중력을 되찾기}

\begin{itemize}
  \item \textbf{핵심 사상}: "스마트폰을 적게 보는 것"이 아니라 "집중할 때 집중하고, 연락할 때 연락하는 것".
  \item \textbf{사용 권장}:\newline · 업무 모드: "별표" 연락처와 필수 앱만 통과 허용;\newline · 수면 모드: 자동 활성화, 알람만 허용;\newline · 운동 / 독서 / 회의: 현재 활동에 맞춰 맞춤 설정.
  \item \textbf{고급 활용}: 자동화 + 위치/시간 트리거를 사용하여 모드가 자동으로 전환되도록 합니다. 예: iOS 단축어에서 "회사에 도착하면 자동으로 업무 모드 진입" 설정 가능.
\end{itemize}

\booksubsection{알림 관리: 빨간 점이 더 이상 “끊임없는 호출”이 되지 않게}

\begin{itemize}
  \item \textbf{마케팅 알림 끄기}: 90\%의 푸시는 사실 필요하지 않음.
  \item \textbf{중요도순 정렬}: 인간 통신류(전화, 문자, 인스턴트 메시지) 상단 고정; 뉴스, 쇼핑, 소셜 플랫폼류는 "무음 푸시"로 설정 가능.
  \item \textbf{그룹별 알림}: 같은 종류 알림 함께 접기, 한 번 잠금 해제 시 수십 개 메시지에 파묻히는 것 방지.
\end{itemize}

\booksubsection{빠른 제스처: 조작이 한 발 빠르게}

\begin{itemize}
  \item \textbf{스크린샷}: 전원 + 볼륨 하(iOS 풀스크린) / 전원 + 볼륨 하(안드로이드 공통).
  \item \textbf{긴 스크린샷 / 스크롤 스크린샷}: 일반 스크린샷 후 "긴 스크린샷" 옵션 클릭.
  \item \textbf{원터치 QR 스캔}: 제어 센터 / 빠른 설정에 "스캔" 아이콘 추가.
  \item \textbf{뒷면 탭(iPhone)}: [설정] > [접근성] > [터치] > [뒷면 탭]에서 "스크린샷" 또는 "잠금 화면" 등 동작 바인딩.
  \item \textbf{음성 웨이크업}: Siri / 샤오아이 / 샤오이 / 샤오부, 한 번 부르면 날씨 조회, 알람 설정 가능.
\end{itemize}

\booksubsection{입력기: 가장 과소평가된 효율 도구}

\begin{itemize}
  \item \textbf{올바른 입력기 선택}:\newline · Sogou / Baidu: 어휘가 크고 신조어 업데이트가 빠름;\newline · iFlytek: 음성 입력 인식률이 가장 높아 회의록 작성에 최적;\newline · Gboard(Google): 깔끔하고, 크로스 플랫폼 동기화가 강하며, 스와이프 입력 지원;\newline · iOS 기본: 프라이버시가 가장 좋고, 시스템과 깊이 통합.
  \item \textbf{필수 스위치}: 스마트 교정, 연상어, 클라우드 어휘 동기화, 위로 스와이프로 숫자 기호 입력.
  \item \textbf{프라이버시 팁}: "문장 연상" 또는 "개인화 추천" 기능 신중 사용, 입력 내용 업로드 가능.
\end{itemize}

\booksubsection{데스크톱과 멀티태스킹: 작업 흐름을 매끄럽게}

\begin{itemize}
  \item \textbf{App 분할 화면 / PIP}: 영상 시청 시 메시지 회신, 인터넷 시 노트 작성.
  \item \textbf{작은 창 / 플로팅 창}: 위챗 / 도우인 "작은 창 모드" 지원, 회신 완료 후 접기.
  \item \textbf{자주 쓰는 App 상단 고정}: Dock 바(화면 하단)에 가장 빈도 높은 4개만 배치, 검색 시간 감소.
\end{itemize}

\begin{tcbinfo}{"업무 시작 첫날" 점검 리스트}
새 스마트폰을 받거나 대규모 업데이트 후, 10분 정도 투자하여 다음 설정을 확인하세요:\newline (1) 잠금 화면 비밀번호 + 생체 인식 (2) 알림 그룹화 + 마케팅 푸시 끄기 (3) 글꼴 크기 / 주사율 / 다크 모드\newline (4) 기본 입력기 / 브라우저 / 메일 (5) 긴급 연락처 / 의료 카드 (6) 백업 설정

\end{tcbinfo}

\booksection{접근성과 배리어프리: 기술의 선행 힘}

“접근성”이라는 이름에 겁먹지 마세요. 이 기능은 누구나 사용할 수 있습니다. 처음에는 장애인을 위해 설계되었지만, 일반인에게도 매우 유용합니다 - 일시적으로 손이 다친 경우, 아기를 안고 있는 경우, 지하철에서 이어폰이 고장나서 잘 안 들리는 경우, 노안으로 작은 글씨가 안 보이는 경우... 모두 도움이 됩니다. 결국 배리어프리 디자인은\textbf{기술을 선하게 만드는}철학을 구현합니다: 누구나 동등하게 디지털 라이프를 즐길 수 있도록 합니다.

\booksubsection{시각 장애 사용자의 “눈” - 스크린 리더 완전 가이드}

스크린 리더는 시각 장애 사용자가 스마트폰을 사용하는 핵심 도구입니다. 화면의\textbf{모든 시각적 콘텐츠를 읽어줍니다}. 사용자는 제스처와 음성으로 스마트폰을 조작합니다. 아래에서 iOS “VoiceOver”와 안드로이드 TalkBack 사용법을 각각 소개합니다.

\begin{tcbinfo}{iOS VoiceOver(VoiceOver)}
\textbf{활성화 경로}: [설정] > [접근성] > [VoiceOver], 열고 켜기/끄기. 측면 키 3회 탭(또는 홈 화면 키)으로 빠른 켜기/끄기 가능.

\textbf{기본 제스처}:

\begin{enumerate}
  \item 한 손가락 스와이프: 화면 내용 픽셀별 탐색;
  \item 두 손가락 쓸기: 상하/좌우 페이지 전체 스크롤;
  \item 한 손가락 더블 탭: 현재 포커스 요소 활성화("탭"과 동일);
  \item 세 손가락 스와이프:전환페이지;
  \item 한 손가락 세 번 탭: 확대 또는 선택(일반 세 번 탭과 다름).
\end{enumerate}

\textbf{개인화 설정}: [VoiceOver] > [음성]에서\textbf{읽기 속도}(기본 약 50\%, 70-80\%로 조절하면 더 매끄럽게 들림),\textbf{음높이}를 조절하거나 다른 음성 팩으로 전환할 수 있습니다. iOS에는 “향상된 음성”과 “음성 컨트롤러”가 내장되어 있어, 전자는 호흡음과 멈춤을 줄여 더 자연스럽게 들립니다.

\end{tcbinfo}

\begin{tcbinfo}{안드로이드 TalkBack}
\textbf{활성화 경로}: [설정] > [접근성] > [TalkBack]에서 스위치를 켭니다. 볼륨 위아래 버튼을 3초간 길게 눌러 빠르게 활성화할 수도 있습니다(일부 브랜드 지원).

\textbf{기본 제스처}:

\begin{enumerate}
  \item 한 손가락 좌우 스와이프: 포커스 요소 전환;
  \item 한 손가락 더블 탭: 요소 활성화;
  \item 두 손가락 상하/좌우 스와이프: 스크롤;
  \item 세 손가락 한 번 탭: 읽기 모드 진입(연속 읽기);
  \item 네 손가락 한 번 탭: 읽기 일시정지/계속.
\end{enumerate}

\textbf{개인 설정}: [TalkBack] > [텍스트 음성 변환 설정],\textbf{말 속도},\textbf{음조}조정 가능,\textbf{음성 엔진}전환 가능(Google 텍스트 음성 변환 엔진 설치 권장, 읽기가 더 자연스러움).

\end{tcbinfo}

\booksubsection{작은 글자를 크게, 흐릿함을 선명하게 - 시각 조정}

\booksubsection{확대 제스처}

\begin{itemize}
  \item \textbf{iOS}: [설정] > [접근성] > [확대/축소]에서 활성화 후\textbf{화면을 세 손가락으로 가볍게 탭}하면 확대됩니다. 기본은 전체 화면 확대이며, 창 확대(부분 확대)도 선택 가능합니다.\textbf{두 손가락으로 벌리고 오므리면 배율 조절(1.5배 \textasciitilde{} 15배)}, 세 손가락으로 두 번 탭하면 현재 확대 비율이 고정됩니다.
  \item \textbf{안드로이드}: [설정] > [접근성] > [확대]에서 세 손가락으로 탭하여 확대하고, 다시 세 손가락으로 탭하여 종료; 또는\textbf{화면을 세 손가락으로 세 번 탭}하여 임시 확대 활성화(최대 8배).
\end{itemize}

\booksubsection{고대비 텍스트 / 굵은 텍스트}

\begin{itemize}
  \item \textbf{iOS}: [설정] > [접근성] > [디스플레이 및 텍스트 크기]에서 "굵은 텍스트", "대비 강화", "투명도 감소" 등의 스위치를 활성화. "스마트 반전"은 밝은 배경에서 텍스트가 검은색으로 유지되어 흰 배경에 흰 글자가 되는 것을 방지합니다.
  \item \textbf{안드로이드}: [설정] > [디스플레이] > [글꼴 크기]에서 "소, 보통, 대, 특대" 네 단계로 조절 가능. 화웨이/샤오미는 추가로 "굵게"와 "고대비 텍스트" 스위치를 제공.
\end{itemize}

\booksubsection{색상 필터(색맹 친화)}

\begin{itemize}
  \item \textbf{iOS}: [설정] > [접근성] > [디스플레이 및 텍스트 크기] > [색상 필터]에서\textbf{빨강/초록/파랑/노랑 네 가지 색맹 모드}와 "회색조", "색상 반전" 옵션을 제공. 시스템에서 색각 테스트를 먼저 한 후 해당 필터를 선택할 것을 권장.
  \item \textbf{안드로이드}: [설정] > [접근성] > [색상 보정]에서\textbf{세 가지 색맹 모드}+ "색상 반전" 스위치.
\end{itemize}

\booksubsection{AI 이미지 설명: AI가 당신의 “눈”이 되어 주세요}

스크린 리더가 이미지를 만날 때, 전통적인 방식은 건너뛰거나 “이미지” 두 글자만 읽는 것이었습니다. 하지만 이제 AI가\textbf{한 문장으로 이미지 내용을 설명}할 수 있어, CAPTCHA 인식, 이미지 읽기, 메뉴 식별에 매우 유용합니다.

\begin{tcbinfo}{각 시스템의 AI 이미지 설명 지원 현황}
\begin{itemize}
  \item \textbf{iOS 15+}: [설정] > [접근성] > [VoiceOver] > [음성 설명]에서 "이미지 설명"을 활성화. 시스템이\textbf{온디바이스 AI 모델}을 호출하여 자동으로 이미지 텍스트 설명을 생성(iPhone 12 이상 필요).
  \item \textbf{신규 TalkBack 앱}: Google TalkBack은 최신 버전에서\textbf{이미지 Q\&A}기능을 통합하여, "이 이미지에 무엇이 있나요?"라고 물으면 AI가 실시간으로 답변(안드로이드 시스템 버전에 엄격히 의존하지 않으며, TalkBack 앱 업데이트만으로 사용 가능).
  \item \textbf{타사 앱}: "Seeing AI"(iOS), "Lookout"(안드로이드), "샤오이 방방"(화웨이)는 모두 우수한 AI 시각 어시스턴트로서, 물품 식별, 문서 읽기, 얼굴 표정 인식 가능.
\end{itemize}

\end{tcbinfo}

\begin{tcbpractice}{사용 시나리오}
한 거리 사진을 보고 “이 사진에 몇 명이 있나요? 그들은 무엇을 하고 있나요?” 물어보세요, 또는 배달 App 열고 AI가 메뉴 이미지의 메뉴 이름 읽어주기 -\textbf{이것이 진정한 “기술이 삶을 더 평등하게”}.

\end{tcbpractice}

\booksubsection{청각 장애 사용자의 “귀” - 실시간 자막 완전 가이드}

실시간 자막(Live Captions)은\textbf{시스템 수준의 음성-텍스트 변환}기능으로, 청각 장애 사용자는 물론 시끄러운 환경, 도서관, 외국어 상황에서 매우 유용합니다.

\begin{tcbinfo}{실시간 자막 활성화 및 최적화}
\textbf{iOS}: [설정] > [접근성] > [실시간 자막], 열고 켜기/끄기. 시스템이 기기의 모든 오디오 감지(전화, 영상, 팟캐스트, FaceTime, 음성 메모), 화면 상단에 자동 자막 표시.

\textbf{안드로이드 12+}:\textbf{볼륨 키}누름 팝업의 “실시간 자막” 아이콘 탭; 또는 [설정] > [접근성] > [실시간 자막]에서 켜기.

\textbf{정확도 최적화 팁}:

\begin{enumerate}
  \item 가능한 한\textbf{조용한 환경}에서 사용, 에어컨, 기계 소음 멀리;
  \item 말하는 사람이\textbf{발음이 명확하고 말 속도가 적당할}때 인식률 가장 높음;
  \item 다자 대화는 정확도를 낮춤, "한 명만 표시" 모드 사용 가능;
  \item 중국어 일반어, 광동어, 영어 등 주류 언어 지원,\textbf{소수 언어 인식률은 비교적 낮음};
  \item 말하는 도중 배경 음악이나 사람 목소리 방해 있으면, 인식이 끊길 수 있음.
\end{enumerate}

\end{tcbinfo}

\booksubsection{통화 시나리오 듀얼 채널: 음성-텍스트 변환 + 텍스트 음성 변환}

사용자가\textbf{듣지도 못하고 말도 할 수 없을}때(갑작스러운 실음, 후부 수술 회복기, 심한 말더듬), 스마트폰은 “말 전달원” 역할도 합니다.

\begin{tcbinfo}{통화 실시간 전환(TTY / RTT)}
\textbf{iOS}: [설정] > [접근성] > [RTT/TTY], 켜기 후 통화 시 화면에 키보드 표시, 상대방의 음성이 실시간으로 문자로 변환,\textbf{당신이 입력한 글자도 음성으로 읽어져}상대방에게 전달.

\textbf{안드로이드}: [설정] > [접근성] > [실시간 전사] + [소리 알림]. “실시간 전사”가 상대방의 말을 실시간으로 문자 표시; “소리 알림”은 음성이 불편할 때 당신의 입력을 음성으로 변환 재생.

\textbf{호환성}:

\begin{itemize}
  \item \textbf{RTT(Real-Time Text)}: 양쪽 모두 지원 필요, 글자를 입력하는 대로 표시, 지연 없음, 주류 통신사 이미 지원;
  \item \textbf{TTY(전신 타자)}: 전용 TTY 기기 필요, 이미 점진적으로 RTT로 대체;
  \item 위챗/QQ의\textbf{음성-텍스트 변환 기능}도 대안: 음성 메시지 길게 누르기 "텍스트 변환".
\end{itemize}

\end{tcbinfo}

\booksubsection{모든 사람을 위한 “작지만 실용적” 기능}

\begin{itemize}
  \item \textbf{뒷면 탭}(iOS 14+): [설정] > [접근성] > [터치] > [뒷면 탭],\textbf{2회/3회 탭}으로 스크린샷, 잠금 화면, 제어 센터 등 동작 트리거 바인딩 가능. 양손으로 택배 받을 때 특히 유용.
  \item \textbf{보조 터치 / 플로팅 볼}: iOS [접근성] > [터치] > [보조 터치], 안드로이드 [접근성] > [플로팅 네비게이션]. 가상 버튼으로 제스처 또는 물리 키 대체, 키 고장 또는 한 손 조작 시 유용.
  \item \textbf{스위치 제어}: 외부 스위치, 단일 탭, 머리 동작으로 스마트폰 제어,\textbf{지체 장애 사용자를 위해 설계}. iOS [접근성] > [스위치 제어], 안드로이드 [접근성] > [스위치 제어].
  \item \textbf{소리 인식}(iOS): [접근성] > [소리 인식], 켜기 후 스마트폰이 초인종, 아기 울음소리, 화재 경보, 사이렌 등 소리 "듣기",\textbf{인식 시 진동+화면 알림}, 청각 장애 사용자와 수면 중 매우 가치 있음.
  \item \textbf{LED 점멸 알림}: iOS [접근성] > [오디오/시각] > [LED 점멸 알림], 안드로이드 [설정] > [알림] > [플래시 알림]. 수신 및 알림 시 플래시 점멸,\textbf{시끄러운 환경에서 무음으로도 인식 가능}.
  \item \textbf{모노 오디오}: [설정] > [접근성] > [오디오] > [모노 오디오], 좌우 채널 병합. 한쪽 귀 청력 손상 사용자가 나머지 소리 정보 잃지 않음.
\end{itemize}

\begin{tcbinfo}{애플의 배리어프리 설계 철학}
애플은 배리어프리 분야에 막대한 투자, 매년 접근성 업데이트 및 추가. 애플 공식 홈페이지에 전용 배리어프리 페이지, 각 항목 기능의 사용 방법 상세 소개. iOS의 배리어프리 기능은 보통 안드로이드보다 더 완벽하고 경험이 더 매끄러움, 이것이 많은 장애 사용자가 애플을 선택하는 중요한 이유.

\end{tcbinfo}

\begin{tcbpractice}{30초 체험}
잠시 여가 시간 찾아, 아래 체험 진행:

\begin{enumerate}
  \item \textbf{VoiceOver / TalkBack}켜기, 눈에 의존하지 말고 세 손가락 슬라이드로 위챗 탐색 시도;
  \item \textbf{실시간 자막}열기, 스마트폰에 대고 말하며 인식 정확도 확인;
  \item \textbf{세 손가락 탭}으로 글자 200\% 확대, 시각 장애 사용자의 일상 체험;
  \item 이전 모든 기능 끄기, "만약 내가 영원히 이 기능들이 필요하다면, 어떻게 설계되길 바랄까" 생각.
\end{enumerate}

\end{tcbpractice}
